\section{Implementation, scope reuse, and validation}

\subsection{The uniqueness checker}
The complete native checker uses a tagged JSON syntax with de Bruijn
indices.  Terms are \texttt{["v",i]}, \texttt{["z"]},
\texttt{["s",t]}, and \texttt{["i",n,a,h]}; only the last child of
the iterator binds a numerical variable.  Formulas are
\texttt{["=",s,t]}, \texttt{["q",B]}, and \texttt{["d",B]}.
For $D$, the closure coordinate is index zero and the witness is index one.
A certificate has exactly the fields \texttt{variables}, \texttt{context},
\texttt{goal}, and \texttt{proof}.  All are checked by the standalone
module, including malformed values and Boolean/integer aliases.

The new proof instruction is
\texttt{["ur",f,h,t,p]}.  It validates $f$ and $h$ at one extra numerical
slot and $t$ at the original scope.  It checks $p$ in the lifted original
context, with no added hypothesis, against
\[
 f(Sz)=h(f(z)).
\]
It returns $f(t)=I(t,f(0);h)$.  The state slot of $h$ is the sole active
slot; the checker does not permit an additional counter slot in the
update.  Its matched control uses exactly the same syntax, traversals,
equality rule, logical rules, JSON interface, and errors, replacing UR
by equation-progress introduction.  The 71-byte increase is therefore a
real cost of this direct rule replacement, not a difference in validation.

The supplied tests exercise correct and incorrect UR recurrences, retention
of the original context, absence of an induction hypothesis, freshness
of the target, nested iterator binders, passive-parameter shifts, and
eigenparameter escape.  The 66 process checks run both normal Python and
\texttt{python -O}; all have their expected outcomes.  These are kernel
regression checks.  They are not a formal verification of the complete
metatheorem or certificates for every arithmetic lemma in this report.

\subsection{The separate four-byte refinement}
The 866-byte program implements the existing scalar calculus, not UR.
It composes two existing implementation observations: lexical validation
of the entire integer/list certificate, and reuse of the first kernel
scope traversal for public-goal bound-variable validation
\cite{validated,scope}.  No numerical axiom or inference rule changes.

The only source change is the following condition in the public grammar
validator:
\begin{center}
\texttt{2 < a < d}\quad becomes\quad\texttt{a > 2}.
\end{center}
The remaining checks still establish the public term/formula arities,
the permitted leaf roles, and integer/list types.  At the initial kernel
call, the existing scope map traverses the entire goal before attempting
any proof instruction.  A nonzero term variable outside its binder depth
reaches the map's invalid-substitution sentinel and rejects.  Thus an
additional goal admitted by the relaxed preliminary test cannot be
accepted by the complete checker.  Valid goals reach the byte-identical
kernel and the same proof stream.  This proves equality of the acceptance
sets on the documented finite JSON interface; the accompanying detailed
argument records the exact sentinel and depth convention.

The improvement is four normalized bytes, two significant tokens, three
AST nodes, and three gzip bytes in this implementation.  Its relevance is
that it is a measured full-program improvement with the same public
contract.  Its mathematical significance is limited: it is not a new
axiomatization, and it does not beat the distinct compact guarded codec.
The differential suite makes 332 standalone runs, equally split between
ordinary and optimized Python, with 108 acceptances, 224 rejections, and
zero mismatches.  It includes malformed JSON, unused-prefix type attacks,
freshness regressions, and nine goals that pass the relaxed preliminary
test but are rejected by the initial kernel scope traversal before any
proof packet is consumed.  The exact code-change check and the scoped
argument about recursion limits are included.

\subsection{Arithmetic validation and assurance limits}
Two independent standard-natural checks accompany the proof.  One checks
the complete explicitly defined scalar predecessor construction and then
the reflection identities using subtraction implemented by iterating that
predecessor.  It reports 1884 successful checks.  The second checks 477760
instances of the Boolean comparisons, accumulator, telescoping equations,
and equality test, with an independent finite-sum oracle for the
accumulator and four deliberately incorrect nearby variants that are
detected.  Bounds and exact counts are recorded in the JSON results.

The universal result rests on the written derivations.  Finite tests
cannot establish it in nonstandard structures or intuitionistic models.
The new equivalence has been reviewed independently at the predecessor,
reflection, and logical-bootstrap boundaries, but is not yet formalized
in Lean or Rocq.  The already existing scalar-to-iPRA bridge is the only
proof-assistant verification inherited from the repository.  The new
checkers are executable artifacts, not proof assistants verifying their
own correctness.

\section{Further research questions}
\begin{enumerate}[leftmargin=*]
\item Formalize scalar UR and the explicit accumulator in the existing Lean
project, preserving the dependency order: cases and logic, UR predecessor,
Boolean diagonal descent, reflection, then general equation induction.
The Boolean descent lemma should isolate much of the new proof.
\item Can the complete UR checker beat a matched induction checker after
an architectural change that genuinely shares computation and uniqueness?
The present direct replacement is larger; that negative measurement should
guide the next attempt.
\item Can the uniqueness rule be weakened to a restricted family of
updates while retaining the predecessor bootstrap and retraction rigidity?
The reflection stage mostly uses constant-recurrence uniqueness, whereas
predecessor and retractions use nontrivial updates.  A justified reduction
would need to cover both uses.
\item Can the explicit scalar accumulator support a more direct
primitive-recursion compiler or a shorter predecessor construction?
Its recurrence is counter-dependent, yet its definition uses only
composition of scalar subtraction and addition.  This distinction may
help avoid premature tuple coding in other bootstraps.
\item Can a smaller complete presentation be established under one fixed
public syntax and validation contract?  Current reported minima span
different languages and input envelopes.  A common contract is necessary
before a meaningful lower-bound or absolute-minimality question can be
posed.
\end{enumerate}

\section{Reproduction and files}
The archive includes this standalone \file{report.tex} and its compiled
\file{report.pdf}; editable proof sections; the two matched uniqueness
checkers; the 866-byte existing-calculus checker and its exact control;
finite arithmetic and checker tests; measurements and provenance; and
the inherited mathematical source files used in the equivalence bridge.
No network access is required for reproduction.

Run \texttt{python3 reproduce.py} to repeat the supplied measurements and
tests.  Run \texttt{python3 reproduce.py --pdf} to rebuild the report with
pdfLaTeX.  To compile the standalone source directly, run
\texttt{pdflatex -interaction=nonstopmode -halt-on-error report.tex} twice.
The archive contains no claimed complete Lean project for the new theorem.

\begin{thebibliography}{9}
\footnotesize\setlength{\itemsep}{3pt}
\bibitem{project} Vladimir Reshetnikov, \emph{ProveIt: Concise presentations
of intuitionistic primitive recursive arithmetic}.  Repository snapshot
\texttt{11be5aabc6b07262046e2553419bc6b05271d178}.
\url{https://github.com/VladimirReshetnikov/ProveIt/tree/11be5aabc6b07262046e2553419bc6b05271d178/Logic/PrimitiveRecursiveArithmetic/ConciseIPRA}.
\bibitem{scope} \emph{Concise iPRA: scope reuse and mathematical limits},
incoming package \file{ipra_scope_kernel_2026-11-11.zip}, especially
\file{sections/uniqueness_note.tex}, in the pinned project snapshot.
The exact inherited uniqueness note is reproduced in \file{reference/}.
\bibitem{choice} \emph{Complementary introductions with one local assumption},
\file{Research/certified_zero/reference/inherited_choice_basis.tex},
in the pinned project snapshot.  Includes the logical reconstruction
independent of progress induction.
\bibitem{validated} \emph{Validated presentation},
\file{Research/validated_presentation/code/native_typed_full.py}
and its validator proof, in the pinned project snapshot.
\bibitem{zero-normal} \emph{Zero-normal scalar iteration}, incoming package
\file{ipra_zero_normal_2026-11-11.zip}, in the pinned project snapshot.
\bibitem{template-free} \emph{Template-free equality transport}, incoming
package \file{ipra_template_free.zip}, in the pinned project snapshot.
\bibitem{okada-scott} M. Okada and P. J. Scott,
\emph{A Note on Rewriting Theory for Uniqueness of Iteration},
Theory and Applications of Categories \textbf{6} (1999), no.~4, 47--64.
\url{https://tac.mta.ca/tac/volumes/6/n4/6-04abs.html}.
\end{thebibliography}
\end{document}
